% القسم: المناهج
% الفصل: الاستقراء
% المبحث: الاستقراء القوي

\documentclass[../../../include/open-logic-section]{subfiles}

\begin{document}

\olfileid{mth}{ind}{str}

\olsection{الاستقراء القوي}

في مبدأ الاستقراء الذي ناقشناه آنفاً، نثبت $P(0)$ ونثبت أيضاً أنه إذا
كان $P(k)$ فإن $P(k+1)$. وفي الجزء الثاني نفترض صدق $P(k)$ ونستعمل هذا
الافتراض لإثبات $P(k+1)$. وبصورة مكافئة، لكل $k>0$، يمكننا افتراض
$P(k-1)$ واستعماله لإثبات $P(k)$؛ والمهم أن نكون قادرين على إجراء
الاستدلال من كل عدد إلى خلفه، أي أن نستطيع إثبات القضية المعنية لكل عدد
على فرض أنها تصح لسابقه.

وثمة صورة أخرى لمبدأ الاستقراء لا نفترض فيها أن القضية تصح لسابق $k$،
وهو $k-1$، فحسب، بل نفترض أنها تصح لجميع الأعداد الأصغر من~$k$، ثم
نستعمل هذا الافتراض لإثبات القضية لـ$k$. وتعطينا هذه الصورة أيضاً
القضية $P(n)$ لكل~$n \in \Nat$. فمتى أثبتنا $P(0)$ نكون بذلك قد
أثبتنا أن $P$ تصح لجميع الأعداد الأصغر من~$1$. وإذا علمنا أنه إذا كانت
$P(l)$ صادقة لكل $l<k$ فإن $P(k)$ صادقة، فنحن نعلم ذلك بوجه خاص عندما
$k=1$. فنستنتج $P(1)$. وبهذا نكون قد أثبتنا $P(0)$ و$P(1)$، أي
$P(l)$ لكل $l<2$؛ وبما أن لدينا أيضاً القضية الشرطية: إذا كانت $P(l)$
صادقة لكل $l<2$ فإن $P(2)$ صادقة، نستنتج~$P(2)$، وهكذا.

والواقع أننا، إذا استطعنا إثبات القضية الشرطية العامة «لكل $k$، إذا
كانت $P(l)$ صادقة لكل $l<k$ فإن $P(k)$ صادقة»، فلن نحتاج إلى إثبات
$P(0)$ على حدة، لأنها تنتج منها. تذكر أن قضية عامة مثل «لكل $l<k$،
$P(l)$» تكون صادقة إذا لم يوجد أي $l<k$. وهذه حالة من التكميم الخالي:
تكون «كل $A$ هي $B$» صادقة إذا لم توجد أي $A$؛ وتكون
$\lforall[x][(!A(x) \lif !B(x))]$ صادقة إذا لم يحقق أي $x$ الشرط
~$!A(x)$. والصيغة الصورية في حالتنا هي
«$\lforall[l][(l < k \lif P(l))]$»؛ وهي صادقة إذا لم يوجد أي
$l < k$. وهذا بالضبط ما يحدث عندما $k=0$: لا يوجد $l<0$، ولذلك تكون
«لكل $l<0$، $P(l)$» صادقة، أياً تكن $P$. ومن ثم يثبت برهان «إذا كانت
$P(l)$ صادقة لكل $l<k$ فإن $P(k)$ صادقة» الحالة~$P(0)$ تلقائياً.

تفيد هذه الصورة عندما لا يمكن جعل إثبات القضية لـ$k$ معتمداً على
القضية لـ$k-1$ وحده، بل قد يتطلب افتراض صدقها لعدد واحد أو أكثر من
الأعداد $l<k$.

\end{document}
